\chapter{Procesy stochastyczne: od rozkładów do trajektorii}
\label{ch:procesy}
\paragraph{Odsyłacze międzytomowe.} Przypomnienie konstrukcji probabilistycznych znajduje się w tomie I, rozdziale~\ref{I-ch:probability}; statystyczny problem danych zależnych omawia tom IV, rozdział~\ref{IV-ch:time-series}.


\section{Dlaczego pojedyncza zmienna losowa już nie wystarcza?}
Rozkład zmiennej $X$ odpowiada na pytanie, jakie wartości może przyjmować
jeden pomiar. Jeżeli mierzymy temperaturę co godzinę, pozycję cząstki
przez kilka sekund albo liczbę klientów w sklepie w ciągu dnia, to
kolejne pomiary są ze sobą powiązane. Sama znajomość rozkładów
$X_0,X_1,\ldots$ nie wystarcza do odpowiedzi na pytania o tę zależność.

Weźmy dwie zmienne o jednakowym rozkładzie Bernoulliego z parametrem
$1/2$. Mogą być niezależne, a mogą być wręcz identyczne: $X_1=X_0$.
W obu przypadkach $\Prob(X_0=1)=\Prob(X_1=1)=1/2$, lecz
$\Prob(X_0=X_1)$ wynosi odpowiednio $1/2$ albo $1$.
Zrozumienie procesu wymaga więc danych o \emph{wspólnych} rozkładach
pomiarów, a potem także informacji o zachowaniu trajektorii.

\begin{definicja}[Proces stochastyczny]
Niech $(\Omega,\calF,\Prob)$ będzie przestrzenią probabilistyczną,
$T\ne\varnothing$ dowolnym zbiorem indeksów, a $(E,\calE)$
przestrzenią mierzalną. \emph{Procesem stochastycznym o przestrzeni
stanów $E$ i zbiorze czasów $T$} nazywamy rodzinę
$X=(X_t)_{t\in T}$, w której każde
$X_t:(\Omega,\calF)\to(E,\calE)$ jest odwzorowaniem mierzalnym.
Dla ustalonego $t$ zmienna $X_t$ opisuje stan w chwili $t$;
dla ustalonego $\omega\in\Omega$ funkcja
$X(\omega):T\to E$, $t\mapsto X_t(\omega)$, jest
\emph{trajektorią} lub \emph{realizacją} procesu.
\end{definicja}

Typowe przypadki to $T=\N_0:=\{0,1,2,\ldots\}$ (czas dyskretny)
oraz $T=[0,\infty)$ (czas ciągły). Przestrzenią stanów może być
$\R^d$, zbiór skończony, przestrzeń funkcji czy przestrzeń miar.
Ważne, że sam proces jest rodziną funkcji mierzalnych; nie
wymagamy jeszcze mierzalności \emph{wspólnej} odwzorowania
$(t,\omega)\mapsto X_t(\omega)$, ani żadnej ciągłości trajektorii.

\begin{center}
\begin{tikzpicture}[>=Latex,scale=0.95]
\node[draw,rounded corners,minimum width=2.9cm,minimum height=2.5cm,fill=blue!3] (a) at (0,0) {};
\node[above=3pt] at (a.north) {losowy wynik $\omega$};
\node at (-0.8,0.35) {$\omega_1$};
\node at (0.6,0.35) {$\omega_2$};
\node at (-0.3,-0.55) {$\omega_3$};
\node[draw,rounded corners,minimum width=5.8cm,minimum height=2.5cm,fill=green!3] (b) at (7.2,0) {};
\node[above=3pt] at (b.north) {trajektoria $t\mapsto X_t(\omega)$};
\draw[->] (4.8,-0.8)--(9.6,-0.8) node[right] {$t$};
\draw[->] (4.85,-0.85)--(4.85,0.8) node[above] {$x$};
\draw[thick,blue!70!black] (5,-0.35)--(5.7,0.1)--(6.5,-0.15)--(7.25,0.58)--(8.05,0.25)--(8.8,0.65)--(9.3,0.3);
\draw[->,thick] (1.7,0.3) -- (4.25,0.3) node[midway,above] {$\omega\mapsto X(\omega)$};
\end{tikzpicture}
\par\smallskip\small Schemat: jedna realizacja losowości wyznacza całą funkcję czasu, nie tylko jeden pomiar.
\end{center}

\begin{przyklad}[Losowa prosta]
Niech $Z$ będzie rzeczywistą zmienną losową i $X_t=tZ$ dla $t\ge0$.
W chwili $t$ rozkład $X_t$ jest obrazem rozkładu $Z$ przez
odwzorowanie $z\mapsto tz$. Każda trajektoria jest prostą o losowym
nachyleniu. Zmienność po czasie jest całkowicie zdeterminowana przez
jedną liczbę $Z(\omega)$: znajomość $X_1$ wyznacza wszystkie $X_t$.
Jeśli $\E Z=0$ i $\Var(Z)=\sigma^2<\infty$, to
$\Cov(X_s,X_t)=st\sigma^2$. Przyrosty
$X_t-X_s=(t-s)Z$ są zwykle silnie zależne; nie tworzą procesu o
niezależnych przyrostach.
\end{przyklad}

\begin{przyklad}[Dwa procesy o tych samych rozkładach pojedynczych]
Niech $\xi_0,\xi_1,\ldots$ będą niezależnymi zmiennymi Bernoulliego
z parametrem $1/2$, a $Z$ kolejną taką zmienną. Definiujemy
$X_n=\xi_n$ i $Y_n=Z$ dla każdego $n\in\N_0$.
Wówczas $\Law(X_n)=\Law(Y_n)$ dla każdego $n$, ale
$\Prob(X_0=X_1)=1/2$, podczas gdy $\Prob(Y_0=Y_1)=1$.
Dla $n\ge1$ mamy również
$\Prob(X_0=\cdots=X_n)=2^{-n}$,
a $\Prob(Y_0=\cdots=Y_n)=1$.
To pierwszy powód, aby badać rozkłady \emph{skończeniewymiarowe}.
\end{przyklad}

\section{Rozkłady skończeniewymiarowe i prawo procesu}
\begin{definicja}[Rozkład skończeniewymiarowy]
Niech $X=(X_t)_{t\in T}$ będzie procesem o stanach w $(E,\calE)$.
Dla każdego $m\ge1$ oraz $t_1,\ldots,t_m\in T$ definiujemy
rozkład wektora $(X_{t_1},\ldots,X_{t_m})$ na
$(E^m,\calE^{\otimes m})$ przez
\[
\mu_{t_1,\ldots,t_m}(A)
:=\Prob\big((X_{t_1},\ldots,X_{t_m})\in A\big),
\qquad A\in\calE^{\otimes m}.
\]
Rodzinę wszystkich miar $\mu_{t_1,\ldots,t_m}$ nazywamy rodziną
\emph{rozkładów skończeniewymiarowych} procesu.
\end{definicja}

Na przykład dla $A_0,\ldots,A_n\in\calE$ mamy
\[
\mu_{t_0,\ldots,t_n}(A_0\times\cdots\times A_n)
=\Prob(X_{t_0}\in A_0,\ldots,X_{t_n}\in A_n).
\]
Przecinki wewnątrz $\Prob(\ldots)$ oznaczają \emph{przecięcie zdarzeń}.
W szczególności $\Prob(X_s\in A,X_t\in B)$ nie jest na ogół
iloczynem $\Prob(X_s\in A)\Prob(X_t\in B)$; taki iloczyn
wymaga niezależności zmiennych $X_s$ i $X_t$.

\begin{twierdzenie}[Zgodność rzeczywistych rozkładów procesu]
Niech $X=(X_t)_{t\in T}$ będzie procesem o stanach w $(E,\calE)$.
Wówczas dla dowolnych $m\ge1$, $t_1,\ldots,t_m\in T$,
permutacji $\pi$ zbioru $\{1,\ldots,m\}$ oraz każdego
$A\in\calE^{\otimes m}$ zachodzi
\[
\mu_{t_{\pi(1)},\ldots,t_{\pi(m)}}(A)
=\mu_{t_1,\ldots,t_m}(\Pi_\pi^{-1}(A)),
\quad
\Pi_\pi(x_1,\ldots,x_m):=(x_{\pi(1)},\ldots,x_{\pi(m)}).
\]
Ponadto dla $k<m$ i $B\in\calE^{\otimes k}$ zachodzi
$\mu_{t_1,\ldots,t_m}(B\times E^{m-k})=\mu_{t_1,\ldots,t_k}(B)$.
Jeśli $t_i=t_j$, rozkład $\mu_{t_1,\ldots,t_m}$ jest skupiony na
zbiorze $\{(x_1,\ldots,x_m):x_i=x_j\}$, o ile diagonalny zbiór jest mierzalny
(np. gdy $E$ jest przestrzenią polską z borelowskim $\sigma$-ciałem).
\end{twierdzenie}
\begin{proof}
Równość dotycząca permutacji wynika wprost z równości zmiennych
wektorowych
$(X_{t_{\pi(1)}},\ldots,X_{t_{\pi(m)}})
=\Pi_\pi(X_{t_1},\ldots,X_{t_m})$.
Własność marginalizacji otrzymujemy z tożsamości zdarzeń
$\{(X_{t_1},\ldots,X_{t_m})\in B\times E^{m-k}\}
=\{(X_{t_1},\ldots,X_{t_k})\in B\}$.
Jeżeli $t_i=t_j$, to $X_{t_i}(\omega)=X_{t_j}(\omega)$
dla każdego $\omega$, więc wektor należy do wskazanej diagonali
z prawdopodobieństwem $1$.
\end{proof}

W drugą stronę powstaje podstawowe pytanie: \emph{jeśli zadamy
rodzinę miar o tych własnościach, czy istnieje proces mający właśnie
te rozkłady?} Odpowiemy na nie w rozdziale~\ref{ch:konstrukcja}.

\begin{definicja}[Prawo procesu na przestrzeni iloczynowej]
Niech $(E,\calE)$ będzie przestrzenią mierzalną, a
$\calE^{\otimes T}$ najmniejszym $\sigma$-ciałem na $E^T$,
które czyni wszystkie rzuty $\pi_t:E^T\to E$, $\pi_t(x)=x(t)$,
mierzalnymi. Dla procesu $X=(X_t)_{t\in T}$ jego
\emph{prawem na przestrzeni iloczynowej} nazywamy
$\Law(X):=\Prob\circ\mathbf X^{-1}$, gdzie
$\mathbf X:\Omega\to E^T$ jest dane wzorem
$\mathbf X(\omega)=(X_t(\omega))_{t\in T}$.
\end{definicja}

\begin{stwierdzenie}[Mierzalność i jednoznaczność prawa]
Dla dowolnego procesu $X$ odwzorowanie
$\mathbf X:(\Omega,\calF)\to(E^T,\calE^{\otimes T})$ jest
mierzalne. Jeśli procesy $X,Y$ o tym samym zbiorze czasów $T$
i przestrzeni stanów $(E,\calE)$ mają jednakowe rozkłady
skończeniewymiarowe, to $\Law(X)=\Law(Y)$ na
$(E^T,\calE^{\otimes T})$.
\end{stwierdzenie}
\begin{proof}
Przeciwobraz zbioru $\pi_t^{-1}(A)$ przez $\mathbf X$ jest
$\{X_t\in A\}\in\calF$. Zbiory te generują $\calE^{\otimes T}$,
więc $\mathbf X$ jest mierzalne. Dla drugiej części rozważmy
klasę cylindrów $\pi_{t_1,\ldots,t_m}^{-1}(A)$,
$A\in\calE^{\otimes m}$. Tworzą one $\pi$-układ, gdyż przecięcie
dwóch takich cylindrów jest cylindrem zależnym od sumy ich
skończonych zbiorów indeksów. Obie miary praw procesów
są równe na cylindrach na mocy założenia o rozkładach
skończeniewymiarowych. Ponieważ cylindry generują
$\calE^{\otimes T}$, twierdzenie o jednoznaczności miar
na generującym $\pi$-układzie kończy dowód.
\end{proof}

\begin{uwaga}[Co jest, a co nie jest ustalone przez rozkłady?]
Rozkłady skończeniewymiarowe wyznaczają prawo procesu na
\emph{cylindrycznym} $\sigma$-ciele $\calE^{\otimes T}$.
Nie znaczy to, że dowolny zbiór trajektorii w $E^T$ jest mierzalny
ani że własność ``trajektoria jest ciągła'' da się badać
wyłącznie przez dowolną, nieprzeliczalną operację przecięcia
zdarzeń. Aby poprawnie mówić o ciągłych trajektoriach, będziemy
konstruować modyfikacje ciągłe i korzystać z mierzalnej przestrzeni
$C([0,T];\R^d)$.
\end{uwaga}


\begin{zadanie}[Prawdopodobieństwo dwóch pomiarów i maksimum trajektorii]
Niech $Z$ ma rozkład $\Normal(0,1)$, a $X_t=tZ$ dla
$t\in[0,2]$. Obliczyć $\Prob(X_1>0,X_2>1)$ oraz
$\Prob(\sup_{0\le t\le2}X_t>1)$ i wyjaśnić, czy w obu pytaniach
wystarczy rozkład pojedynczego pomiaru.

\emph{Rozwiązanie.} Ponieważ $X_1=Z$, $X_2=2Z$,
zdarzenie $\{X_1>0,X_2>1\}$ jest równe $\{Z>1/2\}$.
Oznaczając dystrybuantę standardowego rozkładu normalnego przez
$\Phi$, otrzymujemy $1-\Phi(1/2)$.
Dla ustalonego wyniku $\omega$ funkcja $t\mapsto tZ(\omega)$
rośnie, gdy $Z(\omega)\ge0$, a maleje, gdy $Z(\omega)<0$.
Wobec tego $\sup_{0\le t\le2}X_t=\max(0,2Z)$ i również
$\Prob(\sup_{t\le2}X_t>1)=1-\Phi(1/2)$.
Zgodność liczb jest szczególną cechą tego modelu, a nie ogólnym
prawem: w drugim pytaniu występuje funkcjonał \emph{całej}
trajektorii. Dla procesu o ciągłych trajektoriach maksimum można
też wyrazić jako supremum po przeliczalnym zbiorze chwil
wymiernych, co zapewnia mierzalność zdarzenia.
\end{zadanie}

\section{Modyfikacje i nierozróżnialność}
W czasie dyskretnym można zsumować przeliczalnie wiele zbiorów
zerowych. W czasie ciągłym jest nieprzeliczalnie wiele chwil;
nie wolno zatem automatycznie zamieniać zdań ``dla każdego $t$
prawie na pewno'' oraz ``prawie na pewno dla każdego $t$''.

\begin{definicja}[Modyfikacja i nierozróżnialność]
Niech $X,Y$ będą procesami na tej samej przestrzeni probabilistycznej,
o wspólnym zbiorze indeksów $T$ i przestrzeni stanów $E$.
Mówimy, że $Y$ jest \emph{modyfikacją} (wersją) $X$, jeśli
dla każdego $t\in T$ zdarzenie $\{X_t=Y_t\}$ jest mierzalne
i zachodzi $\Prob(X_t=Y_t)=1$.
Warunek mierzalności jest automatyczny, gdy $E$ jest przestrzenią
metryczną z borelowskim $\sigma$-ciałem: diagonala
$\{(x,y)\in E^2:x=y\}$ jest wtedy zbiorem domkniętym.
Mówimy, że $X$ i $Y$ są \emph{nierozróżnialne}, jeżeli istnieje
$N\in\calF$ z $\Prob(N)=0$ takie, że dla wszystkich
$\omega\notin N$ i wszystkich $t\in T$ zachodzi
$X_t(\omega)=Y_t(\omega)$.
\end{definicja}

\begin{stwierdzenie}[Czas przeliczalny]
Jeśli $T$ jest przeliczalny, a $X$ i $Y$ są swoimi
modyfikacjami, to są nierozróżnialne.
\end{stwierdzenie}
\begin{proof}
Dla każdego $t\in T$ zbiór
$N_t=\{\omega:X_t(\omega)\ne Y_t(\omega)\}$ jest mierzalny
i ma miarę zero (przyjmujemy mierzalną diagonalę w $E\times E$).
Wtedy $N=\bigcup_{t\in T}N_t$ jest przeliczalną sumą
zbiorów zerowych. Poza $N$ zachodzi $X_t=Y_t$ dla
wszystkich $t$ równocześnie.
\end{proof}

\begin{przyklad}[Modyfikacja, która nie jest nierozróżnialna]
Niech $\Omega=[0,1]$ z miarą Lebesgue'a i niech
$U(\omega)=\omega$. Dla $t\in[0,1]$ definiujemy
$X_t(\omega)=0$ oraz $Y_t(\omega)=\ind{\{U(\omega)=t\}}$.
Dla ustalonego $t$ zdarzenie $\{U=t\}$ jest pojedynczym
punktem i ma prawdopodobieństwo $0$, a zatem
$\Prob(X_t=Y_t)=1$. Natomiast dla \emph{każdego}
$\omega\in[0,1]$ w chwili $t=U(\omega)$ mamy
$Y_t(\omega)=1\ne X_t(\omega)$. Nie istnieje więc
nawet jeden wynik $\omega$, dla którego trajektorie byłyby
wszędzie równe. Trajektorie $Y$ mają losowy ``kolec''.
\end{przyklad}

\begin{twierdzenie}[Ciągłe modyfikacje są nierozróżnialne]
Niech $T=[0,T_0]$ dla $T_0>0$, a $E$ będzie przestrzenią
metryczną. Niech procesy $X,Y$ mają trajektorie ciągłe na $T$
poza ewentualnymi, osobnymi mierzalnymi zbiorami zerowymi.
Jeżeli $X$ i $Y$ są modyfikacjami, to są nierozróżnialne.
\end{twierdzenie}
\begin{proof}
Ustalmy przeliczalny zbiór gęsty
$D=(\Q\cap[0,T_0])\cup\{T_0\}$.
Dla każdego $t\in D$ mamy $\Prob(X_t=Y_t)=1$,
więc poza jednym mierzalnym zbiorem zerowym $N_1$
zachodzi $X_t(\omega)=Y_t(\omega)$ dla wszystkich $t\in D$.
Do $N_1$ dodajmy zbiory zerowe, na których którakolwiek z
trajektorii może nie być ciągła, otrzymując $N$.
Dla $\omega\notin N$ i dowolnego $s\in[0,T_0]$
wybierzmy $t_n\in D$ z $t_n\to s$.
Wówczas przez ciągłość obu trajektorii
$X_s(\omega)=\lim_n X_{t_n}(\omega)
=\lim_n Y_{t_n}(\omega)=Y_s(\omega)$.
Zbiór $N$ nie zależy od $s$, co jest kluczowym elementem dowodu.
\end{proof}

\section{Przestrzeń trajektorii ciągłych}
\begin{definicja}[{Przestrzeń $C([0,T];\R^d)$}]
Dla $T>0$ i $d\in\N$ oznaczamy przez
$C([0,T];\R^d)$ zbiór wszystkich funkcji ciągłych
$x:[0,T]\to\R^d$ i wyposażamy go w metrykę
$\norm{x-y}_\infty:=\sup_{t\in[0,T]}\norm{x(t)-y(t)}$.
Jego borelowskie $\sigma$-ciało oznaczamy przez
$\calB(C([0,T];\R^d))$.
\end{definicja}

\begin{twierdzenie}[Mierzalność współrzędnych i generowanie $\sigma$-ciała]
Niech $T>0$, $d\ge1$ oraz
$D=(\Q\cap[0,T])\cup\{T\}$.
Dla każdego $t\in[0,T]$ ewaluacja
$e_t:C([0,T];\R^d)\to\R^d$, $e_t(x)=x(t)$,
jest ciągła. Ponadto
\[
\calB(C([0,T];\R^d))=\sigma(e_t:t\in D).
\]
W konsekwencji, jeśli proces $X=(X_t)_{t\in[0,T]}$
ma ciągłe trajektorie dla każdego $\omega$, to odwzorowanie
$\omega\mapsto(t\mapsto X_t(\omega))$ z $\Omega$ do
$C([0,T];\R^d)$ jest mierzalne.
\end{twierdzenie}
\begin{proof}
Ponieważ $\norm{x(t)-y(t)}\le\norm{x-y}_\infty$,
ewaluacja jest $1$-Lipschitzowska, a zatem ciągła.
Stąd $\sigma(e_t:t\in D)\subseteq\calB(C)$.

Dla dowolnego $f\in C$ i $r>0$, z ciągłości funkcji
$t\mapsto\norm{x(t)-f(t)}$ i gęstości $D$ wynika
$\norm{x-f}_\infty=\sup_{t\in D}\norm{x(t)-f(t)}$.
Otwartą kulę można opisać następująco:
\[
B(f,r)=\bigcup_{k\ge1:\,1/k<r}
\bigcap_{t\in D}
\left\{x:\norm{x(t)-f(t)}\le r-\frac1k\right\}.
\]
Każdy zbiór po prawej stronie jest mierzalny względem
$\sigma(e_t:t\in D)$. Aby przejść od kul do wszystkich
zbiorów otwartych, użyjemy separowalności $C$.
Funkcje łamane, których węzły mają racjonalne chwile
(z dołączeniem $0,T$), a wartości w węzłach należą do
$\Q^d$, tworzą zbiór przeliczalny. Jest on gęsty w
$C$: jednostajna ciągłość pozwala przybliżyć funkcję
odpowiednio drobną interpolacją liniową, a wartości
w skończonej liczbie węzłów przybliżyć liczbami wymiernymi.
Wobec tego $C$ ma przeliczalną bazę otwartych kul,
więc każdy zbiór otwarty jest przeliczalną sumą takich kul.
Dowodzi to odwrotnej inkluzji $\calB(C)\subseteq\sigma(e_t:t\in D)$.
Ostatni wniosek wynika z mierzalności $e_t\circ X=X_t$
dla $t\in D$.
\end{proof}

\begin{uwaga}[Dlaczego wymagamy wspólnego zbioru pełnej miary?]
Jeśli trajektorie $X$ są ciągłe tylko prawie na pewno,
to aby traktować $X$ jako zmienną losową o wartościach w $C$,
wybieramy mierzalny zbiór $\Omega_0$ miary $1$, na którym
ciągłość zachodzi \emph{jednocześnie dla wszystkich $t$},
i poza nim przypisujemy dowolną ustaloną trajektorię ciągłą.
Tak zmieniony proces jest modyfikacją poprzedniego, a jego
odwzorowanie do $C$ jest mierzalne z poprzedniego twierdzenia.
Przypadek $C([0,\infty);\R^d)$ można potraktować przy pomocy
metryki jednostajnej zbieżności na przedziałach zwartych.
\end{uwaga}



\begin{uwaga}[Prawo na przestrzeni trajektorii ciągłych]
Jeśli proces $X$ ma trajektorie ciągłe, traktujemy go jako
zmienną losową o wartościach w $C([0,T];\R^d)$. Wówczas
rozkłady skończeniewymiarowe wyznaczają jego prawo
\emph{jako miarę na przestrzeni trajektorii ciągłych}:
równość miar na cylindrach określonych chwilami wymiernymi
przechodzi na całe borelowskie $\sigma$-ciało, ponieważ
te cylindry je generują. Jest to mocniejszy kontekst niż
kanoniczna przestrzeń $E^T$, gdzie sama ciągłość trajektorii
nie musi odpowiadać mierzalnemu zbiorowi cylindrycznemu.
\end{uwaga}

\section{Trajektorie ze skokami: przestrzeń funkcji c\`adl\`ag}
Między procesem o trajektoriach ciągłych a procesem
zupełnie nieregularnym istnieje ważna klasa pośrednia:
trajektorie mogą mieć skoki, ale nie powinny oscylować
w dowolny sposób bezpośrednio przed każdą chwilą.
Taki warunek będzie podstawowy dla procesów Poissona,
łańcuchów Markowa w czasie ciągłym i procesów Lévy'ego.

\begin{definicja}[Trajektoria c\`adl\`ag]
Niech $(E,d)$ będzie przestrzenią metryczną, a $T>0$.
Funkcja $f:[0,T]\to E$ jest \emph{c\`adl\`ag}, jeżeli:
\begin{enumerate}
\item jest prawostronnie ciągła w każdym $t\in[0,T)$,
tzn. $\lim_{s\downarrow t}f(s)=f(t)$;
\item w każdej chwili $t\in(0,T]$ istnieje
granica lewostronna
$f(t-):=\lim_{s\uparrow t}f(s)\in E$.
\end{enumerate}
Przez $D([0,T];E)$ oznaczamy zbiór wszystkich takich funkcji.
Gdy $E=\R^d$, wielkość $\Delta f(t)=f(t)-f(t-)$
(nieokreślona w $0$) nazywamy \emph{skokiem} w chwili $t$.
\end{definicja}

\begin{przyklad}[Losowy skok, ale uporządkowana trajektoria]
Dla $U\sim\operatorname{U}([0,1])$ proces
$X_t=\ind{\{U\le t\}}$ ma trajektorie c\`adl\`ag.
Przed chwilą $U$ przyjmuje wartość $0$,
a w chwili $U$ i później wartość $1$.
Stąd $X_U=1$, $X_{U-}=0$ oraz $\Delta X(U)=1$.
Jest to funkcja prawostronnie ciągła z jednym skokiem.
Jej trajektorie nie są ciągłe, chociaż -- jak wykażemy
w dalszym przykładzie -- proces jest stochastycznie ciągły.
To rozróżnienie motywuje używanie przestrzeni $D([0,T];E)$
zamiast $C([0,T];E)$ przy procesach skokowych.
\end{przyklad}

\begin{uwaga}[Topologia Skorochoda]
Przy zbieżności procesów skokowych zwykła metryka
jednostajna na $D([0,T];\R^d)$ bywa zbyt silna:
przesunięcie chwili pojedynczego skoku o $1/n$
może nie zmniejszać odległości jednostajnej.
Dlatego w bardziej zaawansowanej teorii wprowadza się
topologię Skorochoda $J_1$, dopuszczającą niewielkie
przesunięcia czasu. Na razie potrzebujemy tylko
samego pojęcia trajektorii c\`adl\`ag.
\end{uwaga}

\section{Przyrosty, stacjonarność i regularność w czasie}
\begin{definicja}[Przyrosty]
Niech $X=(X_t)_{t\ge0}$ będzie procesem o wartościach w $\R^d$.
Mówimy, że ma \emph{niezależne przyrosty}, jeżeli dla
każdych $m\ge1$ i $0\le t_0<t_1<\cdots<t_m$
zmienne $X_{t_1}-X_{t_0},\ldots,X_{t_m}-X_{t_{m-1}}$
są niezależne. Proces ma \emph{stacjonarne przyrosty}, jeśli
dla każdych $0\le s<t$ rozkład $X_t-X_s$ zależy tylko od $t-s$.
Mówimy, że $X$ jest \emph{stochastycznie ciągły w chwili $t$},
jeśli dla każdego $\varepsilon>0$ zachodzi
$\lim_{s\to t}\Prob(\norm{X_s-X_t}>\varepsilon)=0$
(granica jest brana w zbiorze indeksów).
\end{definicja}


\begin{twierdzenie}[Rozkład procesu wyznaczony przez niezależne przyrosty]
Niech $X=(X_t)_{t\ge0}$ będzie rzeczywistym procesem
o niezależnych przyrostach i niech $X_0=0$ prawie na pewno.
Dla $0=t_0<t_1<\cdots<t_m$ oznaczmy
$\Delta_j=X_{t_j}-X_{t_{j-1}}$ i
$\varphi_j(u)=\E e^{iu\Delta_j}$, $u\in\R$.
Wówczas dla dowolnych $u_1,\ldots,u_m\in\R$ zachodzi
\[
\E\exp\left(i\sum_{k=1}^m u_k X_{t_k}\right)
=\prod_{j=1}^m\varphi_j\left(\sum_{k=j}^m u_k\right).
\]
W szczególności rozkłady wszystkich przyrostów
$X_t-X_s$, $0\le s<t$, łącznie z niezależnością
przyrostów i warunkiem $X_0=0$, wyznaczają
wszystkie rozkłady skończeniewymiarowe procesu.
\end{twierdzenie}
\begin{proof}
Dla $k=1,\ldots,m$ mamy
$X_{t_k}=\Delta_1+\cdots+\Delta_k$.
Stąd przez przestawienie dwóch \emph{skończonych} sum
$\sum_{k=1}^m u_k X_{t_k}
=\sum_{j=1}^m(\sum_{k=j}^m u_k)\Delta_j$.
Korzystając z niezależności $\Delta_1,\ldots,\Delta_m$,
możemy rozłożyć wartość oczekiwaną iloczynu
wykładników na iloczyn wartości oczekiwanych,
co daje wzór. Funkcja charakterystyczna wyznacza
jednoznacznie rozkład wektora losowego (twierdzenie
o jednoznaczności funkcji charakterystycznej z tomu I),
a dowolną skończoną rodzinę chwil można uporządkować
i dołączyć chwilę $0$. Dowodzi to ostatniego zdania.
\end{proof}

\begin{przyklad}[Dwa pierwsze kroki spaceru losowego]
Niech $\xi_1,\xi_2$ będą niezależne, przyjmujące
wartości $-1,1$ z prawdopodobieństwem $1/2$,
a $S_0=0$, $S_1=\xi_1$, $S_2=\xi_1+\xi_2$.
Prawdopodobieństwo ścieżki $(0,1,0)$ wynosi
$\Prob(\xi_1=1,\xi_2=-1)=1/4$.
Z kolei $\Prob(S_2=0)=1/2$, bo wynik $0$ otrzymujemy
przy ścieżkach przyrostów $(1,-1)$ albo $(-1,1)$.
Stąd $\Prob(S_1=1\mid S_2=0)=(1/4)/(1/2)=1/2$.
Wniosek: nawet gdy pojedyncze przyrosty są niezależne,
położenia $S_1,S_2$ nie są niezależne; drugie położenie
zawiera informację o pierwszym.
\end{przyklad}

\begin{uwaga}
Stacjonarne przyrosty nie oznaczają stacjonarności samego procesu.
Proces Wienera zaczynający w zerze będzie miał stacjonarne
przyrosty, lecz wariancja $X_t$ będzie rosnąć jak $t$.
Stochastyczna ciągłość mówi jedynie o zbieżności w prawdopodobieństwie
przy zbliżaniu czasu, nie o ciągłości \emph{każdej trajektorii}.
Trzeba też odróżnić warunek dla każdej ustalonej chwili $t$
od jednego zdarzenia pełnej miary, na którym wszystkie trajektorie
mają żądaną regularność.
\end{uwaga}

\begin{definicja}[Stacjonarność ścisła i w drugim rzędzie]
Niech $T=\R$ lub $T=[0,\infty)$, a $X_t\in\R^d$.
Proces jest \emph{ściśle stacjonarny}, jeżeli dla każdego
$m\ge1$, $t_1,\ldots,t_m\in T$ i $h$ takich, że wszystkie
$t_i+h\in T$, zachodzi
$(X_{t_1+h},\ldots,X_{t_m+h})\overset{d}= (X_{t_1},\ldots,X_{t_m})$.
Dla procesu rzeczywistego z $\E X_t^2<\infty$ dla wszystkich $t$
mówimy o \emph{stacjonarności w drugim rzędzie}, gdy
$\E X_t=m$ nie zależy od $t$, a $\Cov(X_s,X_t)=c(t-s)$
dla pewnej funkcji $c$ i wszystkich $s,t\in T$.
\end{definicja}

\begin{twierdzenie}
Każdy ściśle stacjonarny proces rzeczywisty,
dla którego $\E X_t^2<\infty$ dla wszystkich $t$,
jest stacjonarny w drugim rzędzie.
\end{twierdzenie}
\begin{proof}
Równość rozkładów $X_t$ i $X_{t+h}$ daje równość średnich
oraz drugich momentów. Równość rozkładów par
$(X_s,X_t)$ i $(X_{s+h},X_{t+h})$ daje równość
$\E[X_sX_t]$ i $\E[X_{s+h}X_{t+h}]$, więc i kowariancji.
Wybierając przesunięcie, które przenosi jedną z chwil na $0$
(gdy $T=[0,\infty)$, używamy dozwolonego przesunięcia
obu nieujemnych chwil), otrzymujemy zależność kowariancji
wyłącznie od różnicy czasów.
\end{proof}

\begin{przyklad}[Stacjonarny proces z gładkimi trajektoriami]
Niech $Z_1,Z_2$ będą niezależnymi zmiennymi standardowo normalnymi
i niech $X_t=Z_1\cos t+Z_2\sin t$ dla $t\in\R$.
Mamy $\E X_t=0$ i
\[
\Cov(X_s,X_t)=\cos s\cos t+\sin s\sin t=\cos(t-s).
\]
Dla dowolnego $m$ wektor $(X_{t_1},\ldots,X_{t_m})$
jest gaussowski, więc wyznaczają go średnie i macierz kowariancji.
Po przesunięciu wszystkich chwil o $h$ te dane nie zmieniają się,
zatem proces jest ściśle stacjonarny. Każda trajektoria jest przy
tym funkcją gładką; stacjonarność nie oznacza ``poszarpania'' ścieżek.
\end{przyklad}

\begin{przyklad}[Ciągłość stochastyczna bez ciągłych trajektorii]
Niech $U$ ma rozkład jednostajny na $[0,1]$ i
$X_t=\ind{\{U\le t\}}$ dla $t\in[0,1]$.
Dla ustalonego $t$ i $s$ zachodzi
$\Prob(\abs{X_s-X_t}>0)=\Prob(U\in(\min(s,t),\max(s,t)])
=\abs{s-t}$, zatem $X$ jest stochastycznie ciągły.
Niemniej dla każdego $\omega$ z $U(\omega)\in(0,1)$
trajektoria ma skok z $0$ do $1$ w chwili $U(\omega)$.
Co więcej, nie istnieje ciągła modyfikacja: gdyby $Y$ była ciągłą
modyfikacją, to poza jednym zbiorem zerowym mielibyśmy
$Y_q=X_q$ dla wszystkich $q\in\Q\cap[0,1]$ jednocześnie.
Dla każdego $\omega$ z tego zbioru pełnej miary oraz
$U(\omega)\in(0,1)$ wartości $X_q(\omega)$ na wymiernych
po obu stronach $U(\omega)$ miałyby różne granice jednostronne.
Ciągła funkcja $Y(\omega)$ nie może mieć takiego zachowania.
\end{przyklad}

\section{Procesy gaussowskie jako zapowiedź konstrukcji}
\begin{definicja}[Proces gaussowski]
Proces rzeczywisty $X=(X_t)_{t\in T}$ nazywamy
\emph{gaussowskim}, jeżeli dla każdego $m\ge1$
i $t_1,\ldots,t_m\in T$ wektor
$(X_{t_1},\ldots,X_{t_m})$ ma wielowymiarowy rozkład
normalny (dopuszczamy rozkłady zdegenerowane).
Jego funkcjami średniej i kowariancji są odpowiednio
$m(t)=\E X_t$ i $K(s,t)=\Cov(X_s,X_t)$.
\end{definicja}

\begin{stwierdzenie}[Dodatnia półokreśloność kowariancji]
Jeśli $X=(X_t)_{t\in T}$ jest procesem rzeczywistym o
skończonych drugich momentach, to funkcja
$K(s,t)=\Cov(X_s,X_t)$ jest symetryczna i dla każdych
$m\ge1$, $t_1,\ldots,t_m\in T$, $a_1,\ldots,a_m\in\R$
zachodzi
$\sum_{i,j=1}^m a_i a_j K(t_i,t_j)\ge0$.
\end{stwierdzenie}
\begin{proof}
Symetria wynika z symetrii kowariancji. Po rozwinięciu
wariancji kombinacji liniowej otrzymujemy
\[
\sum_{i,j=1}^m a_i a_j K(t_i,t_j)
=\Var\left(\sum_{i=1}^m a_i X_{t_i}\right)\ge0.
\]
Każda macierz $[K(t_i,t_j)]_{i,j}$ musi więc być dodatnio
półokreślona, co stanowi istotne ograniczenie przy próbie
zadania procesu przez samą funkcję kowariancji.
\end{proof}

\begin{zadanie}[Czy dwa procesy mają ten sam rozkład?]
Niech $U,V$ będą niezależnymi zmiennymi Bernoulliego z parametrem
$1/2$. Na przedziale $[0,1]$ definiujemy
$X_t=U$ dla wszystkich $t$ oraz
$Y_t=U$ dla $t<1/2$ i $Y_t=V$ dla $t\ge1/2$.
Obliczymy rozkłady jednowymiarowe, prawdopodobieństwo
$\Prob(X_{1/4}=X_{3/4})$ oraz jego odpowiednik dla $Y$.

Dla każdego $t$ obie zmienne $X_t$ i $Y_t$ mają rozkład
Bernoulliego z parametrem $1/2$, więc same rozkłady
jednowymiarowe nie rozróżniają tych procesów. Natomiast
$X_{1/4}=X_{3/4}=U$, zatem
$\Prob(X_{1/4}=X_{3/4})=1$.
Dla drugiego procesu
$Y_{1/4}=U$, $Y_{3/4}=V$; stąd przez niezależność
\[
\Prob(Y_{1/4}=Y_{3/4})
=\Prob(U=V=0)+\Prob(U=V=1)=\tfrac14+\tfrac14=\tfrac12.
\]
Ponadto $X$ ma ciągłe trajektorie, a $Y$ ma skok w chwili
$1/2$ z prawdopodobieństwem $\Prob(U\ne V)=1/2$.
Zadanie pokazuje równocześnie rolę rozkładów łącznych
i regularności trajektorii.
\end{zadanie}

\begin{uwaga}[Granice pierwszego rozdziału]
Dotąd startowaliśmy od \emph{istniejącego} procesu i wyznaczaliśmy
jego prawo. Jeśli chcemy w modelu narzucić rozkłady, np.
$(X_{t_1},\ldots,X_{t_m})\sim\Normal(0,[\min(t_i,t_j)]_{i,j})$,
musimy wykazać istnienie procesu o takich rozkładach.
Dalej będziemy jeszcze wymagać, by jego trajektorie były ciągłe.
Te dwa kroki nie mogą być zlane w jeden.
\end{uwaga}


\begin{przyklad}[Ta sama kowariancja nie oznacza tego samego prawa]
Niech $Z_1,Z_2$ będą niezależnymi zmiennymi standardowo
normalnymi, a $\xi_1,\xi_2$ niezależnymi symetrycznymi
zmiennymi o wartościach $-1,1$. Definiujemy na $\R$
\[
X_t=Z_1\cos t+Z_2\sin t,
\qquad Y_t=\xi_1\cos t+\xi_2\sin t.
\]
Oba procesy mają średnią zero i identyczną funkcję
kowariancji $K(s,t)=\cos(t-s)$, ponieważ
$\E Z_i^2=\E\xi_i^2=1$ oraz wszystkie mieszane
kowariancje są zerowe. Jednak $X_0\sim\Normal(0,1)$,
a $Y_0=\xi_1$ przyjmuje tylko wartości $-1,1$.
Zatem procesy mają inne rozkłady już w chwili $0$.
Funkcja kowariancji wyznacza rozkłady skończeniewymiarowe
\emph{w klasie procesów gaussowskich}, ale nie w klasie
wszystkich procesów o skończonych drugich momentach.
\end{przyklad}

